\documentclass{article}

\usepackage[a4paper, left=1.5cm, right=1.5cm, top=2cm, bottom=2cm]{geometry}

\usepackage{../../../../components/components}

\usepackage{fancyhdr}


% Configuration des en-têtes et pieds de page
\pagestyle{fancy}
\fancyhf{} % reset tout

\fancyhead[L]{DL2 Math-Info PR4}
\fancyhead[C]{Probabilités}
\fancyhead[R]{2025-2026}

\fancyfoot[L]{Ewen Rodrigues de Oliveira}
\fancyfoot[R]{\thepage}

\begin{document}

\docTitle{Chapitre 7 : Outils probabilistes}

\section{Inégalités sur l'espérance}
\subsection{Inégalité de Markov}

\theorem{Proposition}{Inégalité de Markov}{false}{
    Si $X \geq 0$ et $a > 0$, alors :
    \[ \mathbb{P}(X \geq a) \leq \frac{\mathbb{E}[X]}{a} .\]
}

\example{
    On jette $n$ fois une pièce équilibrée et on note $S_n$ le nombre de Piles obtenues.\\
    On a $\mathbb{E}[S_n] = \frac{n}{2}$ par linéarité de l'espérance.\\
    En appliquant l'inégalité de Markov, on obtient :
    \[ \mathbb{P}\left(S_n \geq \frac{3n}{4}\right) \leq \frac{\mathbb{E}[S_n]}{\frac{3n}{4}} = \frac{\frac{n}{2}}{\frac{3n}{4}} = \frac{2}{3} .\]
    L'inégalité est correcte, mais quelque peu décevante car la borne ne décroit pas avec n.
}

\subsection{Inégalité de Bienaymé-Tchebychev}

\theorem{Proposition}{Inégalité de Bienaymé-Tchebychev}{false}{
    Soit $X$ une variable aléatoire réelle discrète de carré intégrable.
    \[
       \forall a > 0, \quad \mathbb{P}(|X - \mathbb{E}[X]| \geq a) \leq \frac{\mathrm{Var}(X)}{a^2} .
    \]
    \textit{En particulier,} si $Var(X) > 0$, on a : $\mathbb{P}(|X - \mathbb{E}[X]| \geq c \sqrt{\mathrm{Var}(X)}) \leq \frac{1}{c^2}$ pour tout $c > 0$.
}

\noindent{\textbf{Preuve :} Appliquer l'inégalité de Markov à la variable aléatoire $(X - \mathbb{E}[X])^2$ qui est positive.}

\remark{L'inégalité de Bienaymé-Tchebychev permet de majorer la probabilité que $X$ s'écarte de sa moyenne de plus de $a$ en fonction de la variance de $X$.}
\example{Reprenons l'exemple du lancer de pièce. On a $Var(S_n) = \frac{n}{4}$. On observe que $\{ S_n \geq \frac{3n}{4} \} \subset \{ |S_n - \mathbb{E}[S_n]| \geq \frac{n}{4} \}$.\\
Puisque $S_n$ est de carré intégrable, on peut appliquer l'inégalité de Bienaymé-Tchebychev :
\[
    \mathbb{P}\left(S_n \geq \frac{3n}{4}\right) \leq \mathbb{P}\left(|S_n - \mathbb{E}[S_n]| \geq \frac{n}{4}\right) \leq \frac{\mathrm{Var}(S_n)}{\left(\frac{n}{4}\right)^2} = \frac{\frac{n}{4}}{\frac{n^2}{16}} = \frac{4}{n} .
\]
Cette majoration est plus satisfaisante que celle obtenue avec l'inégalité de Markov, car elle tend vers 0 lorsque $n$ tend vers l'infini.
}

\subsection{Inégalité de Cauchy-Schwarz}

\theorem{Proposition}{Inégalité de Cauchy-Schwarz}{false}{
    Soit $X$ et $Y$ deux variables aléatoires réelles discrètes de carré intégrable. Alors la variable aléatoire $XY$ est intégrable et on a :
    \[
        |\mathbb{E}[XY]| \leq \sqrt{\mathbb{E}[X^2]} \sqrt{\mathbb{E}[Y^2]} .
    \]
}

\theorem{Corollaire}{Borne sur le coefficient de corrélation}{false}{
    Soit $X$ et $Y$ deux variables aléatoires réelles discrètes de carré intégrable. Le coefficient de corrélation $\rho(X,Y) = \frac{\mathrm{Cov}(X,Y)}{\sqrt{\mathrm{Var}(X)} \sqrt{\mathrm{Var}(Y)}}$ est bien défini et vérifie $|\rho(X,Y)| \leq 1$.
}
\noindent{\textbf{Preuve :} On peut faire une belle preuve en posant dans $\mathbb{R}^2$ le produit scalaire $\langle X, Y \rangle = \mathbb{E}[XY]$ et la norme $\|X\| = \sqrt{\mathbb{E}[X^2]}$. L'inégalité de Cauchy-Schwarz est alors une conséquence directe de l'inégalité de Cauchy-Schwarz dans les espaces euclidiens.}

\remark{L'inégalité de Cauchy-Schwarz est un outil fondamental dans les espaces euclidiens et trouve de nombreuses applications en probabilité, notamment dans l'étude de la covariance et du coefficient de corrélation.}

\newpage
\tableofcontents

\end{document}